% Part: sets-functions-relations
% Chapter: relations
% Section: equivalence-relations

\documentclass[../../../include/open-logic-section]{subfiles}

\begin{document}

\olfileid{sfr}{rel}{eqv}

\olsection{等价关系}

集合上的恒等关系是自反的、对称的和传递的。同时具备这三种性质的关系~$R$ 非常常见。

\begin{defn}[等价关系]
一个自反、对称且传递的关系 $R \subseteq A^2$ 称为\emph{等价关系}。若~$Rxy$，则称 $A$ 的元素 $x$ 和 $y$ 是\emph{$R$-等价的}。
\end{defn}

等价关系引出了\emph{等价类}的概念。一个等价关系将论域“分块”成不同的划分。在每个划分内部，所有对象彼此相关；而来自不同划分的对象则不相关。有时，直接讨论这些划分会很有用。为此，我们引入以下定义：

\begin{defn}\ollabel{def:equivalenceclass}
设 $R \subseteq A^2$ 是一个等价关系。对每个 $x \in A$，$x$ 在~$A$ 中的\emph{等价类}是集合 $\equivrep{x}{R}
= \Setabs{y \in A}{Rxy}$。$A$ 在~$R$ 下的\emph{商集}是 $\equivclass{A}{R} = \Setabs{\equivrep{x}{R}}{x \in A}$，即这些等价类构成的集合。
\end{defn}

下面的命题通过证明等价类确实构成了~$A$ 的划分，从而验证了等价类定义的合理性：

\begin{prop}
如果 $R \subseteq A^2$ 是等价关系，那么 $Rxy$ 当且仅当 $\equivrep{x}{R} = \equivrep{y}{R}$。
\end{prop}

\begin{proof}
从左到右方向：假设 $Rxy$，并设 $z \in \equivrep{x}{R}$。根据定义，有 $Rxz$。由于 $R$ 是等价关系，可得 $Ryz$。（具体而言：由 $Rxy$ 及 $R$ 的对称性得 $Ryx$，再由 $Rxz$ 及 $R$ 的传递性得 $Ryz$。）因此 $z \in \equivrep{y}{R}$。一般化地，有 $\equivrep{x}{R} \subseteq \equivrep{y}{R}$。但完全类似地，有 $\equivrep{y}{R} \subseteq \equivrep{x}{R}$。于是根据外延性原理，$\equivrep{x}{R} =
\equivrep{y}{R}$。

从右到左方向：假设 $\equivrep{x}{R} =
\equivrep{y}{R}$。由于 $R$ 是自反的，$Ryy$，所以 $y \in \equivrep{y}{R}$。于是由假设 $\equivrep{x}{R} = \equivrep{y}{R}$ 也得 $y \in \equivrep{x}{R}$。因此 $Rxy$。
\end{proof}

\begin{ex}
等价关系的一个绝佳例子来自模算术。对任意 $a$, $b$ 和 $n \in \PosInt$，规定 $a \equiv_n b$ 当且仅当 $a$ 除以~$n$ 所得的余数与 $b$ 除以~$n$ 所得的余数相同。（用更形式化的符号表示：$a \equiv_n b$ 当且仅当存在 $k \in \Int$ 使得 $a - b = kn$。）可以看到，对任意~$n$，$\equiv_n$ 都是一个等价关系。并且由~$\equiv_n$ 生成的不同的等价类恰好有 $n$ 个；也就是说，$\equivclass{\Nat}{\equiv_n}$ 有 $n$ 个元素。它们分别是：能被 $n$ 整除的数集，即 $\equivrep{0}{\equiv_n}$；被 $n$ 除余数为~$1$ 的数集，即 $\equivrep{1}{\equiv_n}$；……；以及被~$n$ 除余数为~$n-1$ 的数集，即~$\equivrep{n-1}{\equiv_n}$。
\end{ex}

\begin{prob}
证明对任意 $n \in \PosInt$，$\equiv_n$ 是等价关系，并且 $\equivclass{\Nat}{\equiv_n}$ 恰好有 $n$ 个成员。
\end{prob}

\end{document}